\section{Proofs for geometry and the oracle inequality}
\label{app:oracle-proofs}

\subsection{Elementary transfer and access results}

\begin{proof}[Proof of Proposition~\ref{prop:pinsker-transfer}]
Condition on \((X,R,E=e)\). The difference between expectations of a function with range length \(B_e\) under two probability measures is at most \(B_e\) times their total variation distance. Pinsker's inequality gives
\[
\TV(p_W,q_0)\le\sqrt{\KL(q_0\|p_W)/2}.
\]
Average over \((X,R)\), then use Jensen's inequality for the concave square root.
\end{proof}

\begin{proof}[Proof of Proposition~\ref{prop:access}]
For (i), returning the directly provided feasible parameter makes the student law equal
to the teacher law, so its population divergence is zero.

For (ii), strict monotonicity of the logistic link gives
\(\operatorname{logit}\{q_0(1\mid X_i)\}=X_i^\top\theta_0\).  A full-column-rank
design has the unique solution
\(\theta_0=(X^\top X)^{-1}X^\top z\), where
\(z_i=\operatorname{logit}\{q_0(1\mid X_i)\}\).  The recovered conditional law
therefore has zero population divergence.

For (iii), conditional cross-entropy decomposes pointwise as
\[
\bar\ell_\theta(x,e)=H\{q_0(\cdot\mid x,e)\}
+\KL\{q_0(\cdot\mid x,e)\|p_\theta(\cdot\mid x,e)\}.
\]
At correct specification, its gradient vanishes at \(\theta_0\) for every
prompt.  Under misspecification, the first-order condition and a Taylor
expansion give
\[
0=P_n\nabla\bar\ell_{\wh\theta}
=P_n\nabla\bar\ell_{\theta^*}
+H(\wh\theta-\theta^*)+o_p(n^{-1/2}).
\]
The multivariate central limit theorem yields the displayed sandwich limit. A second Taylor expansion of population risk gives
\[
\bar\cD_w(\wh\theta)-\bar\cD_w(\theta^*)
=\tfrac12(\wh\theta-\theta^*)^\top H(\wh\theta-\theta^*)
+o_p(n^{-1}).
\]
Taking expectations under uniform integrability and using cyclicity of trace yields
\[
\tfrac12\tr\{H(H^{-1}V_XH^{-1})\}
=\tfrac12\tr(V_XH^{-1}).
\]
At correct specification the pointwise gradient is zero, so \(V_X=0\). This is the standard smooth M-estimator calculation~\citep{vandervaart1998}.
\end{proof}

\subsection{TT geometry}

\begin{proof}[Proof of Lemma~\ref{lem:tt-dimension}]
For one layer, the product core space has dimension
\(m_\ell=\sum_jr_{\ell,j-1}m_{\ell j}r_{\ell j}\). Minimal TT representations of the same exact-rank tensor differ by the adjacent invertible transformations displayed in \cref{sec:geometry}; conversely these transformations preserve the contraction. On the minimal-rank stratum the action is free, so each orbit has dimension
\(\sum_{j<q_\ell}\dim\mathrm{GL}(r_{\ell j})
=\sum_{j<q_\ell}r_{\ell j}^2\). The quotient-manifold theorem therefore gives dimension \(m_\ell-\sum_{j<q_\ell}r_{\ell j}^2\)~\citep{holtz2012manifolds}. Independent layer products add dimensions, as do identifiable free coordinates. Tying imposes constraints across this product, so neither the product action nor the additive calculation need remain valid.
\end{proof}

\begin{proof}[Proof of Proposition~\ref{prop:tt-atlas}]
For any orthonormal basis \(U_{\ell j}\) of an unfolding column space,
Cauchy--Binet gives
\[
 1=\det(U_{\ell j}^\top U_{\ell j})
 =\sum_{|I|=r_{\ell j}}\det\{U_{\ell j}[I,:]\}^2.
\]
Consequently at least one row set \(I_{\ell j}\) has determinant in absolute
value at least \(c_{\ell j}\). Because every singular value of the square
submatrix is at most one, its smallest singular value is at least the absolute
determinant. Thus, for at least one pivot, the corresponding principal block
of the projector has smallest eigenvalue at least \(c_{\ell j}^2\). Every
tensor therefore belongs to the corresponding open patch with threshold
\(c_{\ell j}^2/4\) in \cref{eq:pivot-patch}. Counting pivot collections gives
\cref{eq:pivot-count}, including the separate atlas for constrained free
coordinates.

Write \(A=U_{\ell j}[I,:]\). On a thresholded patch, \(A\) is invertible and
\[
 P_{\ell j}[:,I]\{P_{\ell j}[I,I]\}^{-1}
 =U_{\ell j}A^\top(AA^\top)^{-1}
 =U_{\ell j}A^{-1}.
\]
Thus the normalized interface is expressed entirely through the column-space
projector and is independent of every choice of singular vectors. It fixes
the corresponding bond gauge uniquely, and \(\|A^{-1}\|_{\mathrm{op}}<
2/c_{\ell j}\) on the thresholded patch.

On the exact-minimal-rank core space, the derivative kernel of the full
contraction map is the vertical gauge tangent~\citep{holtz2012manifolds}.
Its restriction to the transverse normalized-pivot slice consequently has
zero kernel. Lemma~\ref{lem:tt-dimension} and the constant-rank theorem give
a \(d_\lambda\)-dimensional Euclidean chart around every point of a pivot
patch; they do not assert that the entire patch is one chart.

For uniform control, each unfolding has fixed rank and smallest nonzero
singular value at least \(\eta\). Its column-space spectral projector
\(P_{\ell j}\) is therefore \(C^2\) locally in the tensor, with derivatives
bounded in terms of \(\eta^{-1}\).
Sequential core recovery on a thresholded patch uses contractions, these
spectral projectors, pivot inverses bounded by \(2/c_{\ell j}\), and inverses
of nonzero unfolding blocks bounded by \(\eta^{-1}\). On the compact
\(R\)-bounded set, the quantitative constant-rank theorem therefore supplies
uniform local radii and the derivative bound in
\cref{eq:atlas-constant}. Compactness extracts finitely many such subcharts
over the finite pivot cover; their total number is \(H_\lambda\).

Finally, each subchart is \(C_{\mathrm{atlas},\lambda}\)-Lipschitz. Map the
standard Euclidean \(\epsilon/C_{\mathrm{atlas},\lambda}\)-net of a
\(d_\lambda\)-ball through the chart and use the volume bound
\(N(\delta,B_d(a,t),\norm{\cdot}_2)\le(1+2t/\delta)^{d_\lambda}\).
\end{proof}

\begin{proof}[Proof of Proposition~\ref{prop:block-curvature}]
Let \(S(p)=\operatorname{diag}(p)-pp^\top\). The logit derivative in direction \(H\) is
\[
D_Wz_W[H]=C\operatorname{diag}\{\phi'(Wh)\}Hh.
\]
The conditional softmax information quadratic form is
\[
\{D_Wz_W[H]\}^\top S(p_W)D_Wz_W[H].
\]
The assumed head and activation bounds make this at least
\(\kappa_C\kappa_\phi^2\norm{Hh}_2^2\). Averaging and using
\(\E(hh^\top)\succeq\kappa_hI\) gives
\[
\E\norm{Hh}_2^2=\tr\{H\E(hh^\top)H^\top\}
\ge\kappa_h\norm H_F^2.
\]
The result follows. The head condition also removes the softmax all-ones null direction.
\end{proof}

\subsection{Localized entropy and within-model fitting}

\begin{proof}[Proof of Lemma~\ref{lem:entropy}]
First suppose \(r<g_\lambda\). The global gap puts the entire risk sublevel inside \(U_\lambda^*\), and the margin places it in a Euclidean chart ball of radius \(\sqrt{2r/\mu_\lambda}\) about the projection representative. The pointwise likelihood Lipschitz condition maps a Euclidean \(\delta\)-net to an \(L_\infty\) net of radius at most \(L_\lambda\delta\). The standard volume cover of a \(d_\lambda\)-ball gives
\[
N(u\sqrt r,\mathcal F_\lambda(r),L_\infty)
\le
\left(1+\frac{2L_\lambda\sqrt{2/\mu_\lambda}}{u}\right)^{d_\lambda}.
\]
For \(r\ge g_\lambda\), cover the full compact atlas. If its largest coordinate diameter is \(R_\lambda\), the same argument gives
\[
N(u\sqrt r,\mathcal F_\lambda(r),L_\infty)
\le H_\lambda
\left(1+\frac{2L_\lambda R_\lambda}{u\sqrt{g_\lambda}}\right)^{d_\lambda}.
\]
These displays imply \cref{eq:local-entropy} after enlarging \(a_\lambda\).

Inside \(U_\lambda^*\), pointwise Lipschitzness and the margin give
\[
Q_0^w(\ell_W-\ell_{W_\lambda^*})^2
\le L_\lambda^2\norm{\theta-\theta_\lambda^*}^2
\le\frac{2L_\lambda^2}{\mu_\lambda}
\{\cD_w(W,W^0)-A_w^2(\lambda)\}.
\]
Outside that neighborhood, the likelihood difference is bounded by \(2B\) and the excess risk is at least \(g_\lambda\), which gives the factor \(4B^2/g_\lambda\). Compactness and continuity permit a finite enlargement of the constant on the interface.
\end{proof}

We record the empirical-process step used below. It is a localized finite-dimensional likelihood lemma, stated with the constants needed here.

\begin{lemma}[Localized likelihood deviation]
\label{lem:local-deviation}
Let \(\mathcal F\) be a pointwise measurable, separable, bounded loss-difference class with nonnegative means, variance bound \(Pf^2\le vPf\), and local uniform covering numbers
\[
\log N(u\sqrt r,\{f:Pf\le r\},L_\infty)
\le d\log(a/u)+h.
\]
There is \(C=C(a,v,\norm{\mathcal F}_\infty)\) such that, with probability at least \(1-e^{-z}\), simultaneously over \(f\in\mathcal F\),
\begin{equation}
|(P_n-P)f|\le\tfrac14Pf+C\frac{d+h+z}{n}.
\label{eq:isomorphic-lemma}
\end{equation}
The same conclusion holds for a finite atlas with the \(h=\log H\) term above.
\end{lemma}

\begin{proof}
We give the uniform-metric chaining step because an ordinary
\(L_2(P)\)-cover would not suffice. Fix \(r>0\), put
\(\mathcal F(r)=\{f:Pf\le r\}\), and take proper finite
\(L_\infty\)-nets \(T_k\subset\mathcal F(r)\), not necessarily nested, at radii
\(\epsilon_k=2^{-k}\sqrt r\), \(k\ge0\). Their logarithmic cardinalities are
at most
\[
 d\log(a2^k)+h.
\]
For a net point in \(T_0\), the variance is at most \(vr\) and its envelope
is at most \(\norm{\mathcal F}_\infty\). Bernstein's inequality and a union
bound control every such coarse point by
\[
 C\left\{\sqrt{\frac{vr(d+h+z)}n}
 +\frac{d+h+z}{n}\right\},
\]
where dependence on \(a\) and the envelope is absorbed into \(C\).

Let \(\pi_kf\) be a nearest point to \(f\) in \(T_k\), and connect these
approximants at adjacent levels. Every increment
\(\pi_kf-\pi_{k-1}f\) has envelope at most
\(3\epsilon_{k-1}\) and variance at most
\(9\epsilon_{k-1}^2\). Apply Bernstein simultaneously to all adjacent net
pairs, assigning level \(k\) failure probability proportional to
\(e^{-z-k}\). Summing the resulting bounds uses
\[
 \sum_{k\ge1}2^{-k}\sqrt{k+u}\le C\sqrt{1+u},
 \qquad
 \sum_{k\ge1}2^{-k}(k+u)\le C(1+u).
\]
Separability lets \(\pi_kf\to f\) uniformly. Consequently, with probability
at least \(1-e^{-z}\),
\begin{equation}
 \sup_{f\in\mathcal F(r)}|(P_n-P)f|
 \le C\left\{\sqrt{\frac{vr(d+h+z)}n}
 +\frac{d+h+z}{n}\right\}.
 \label{eq:local-chain-bound}
\end{equation}
The constant also absorbs the bounded uniform diameter in the linear
increment terms.

Finally apply \cref{eq:local-chain-bound} on the dyadic shells
\(2^{j-1}r_0<Pf\le2^jr_0\), assign shell \(j\) failure probability
proportional to \(e^{-z-j}\), and take
\(r_0=C_0(d+h+z)/n\). Young's inequality, together with
\(j+1\le2^{j+1}\), makes the square-root and shell-allocation terms at most
\(Pf/4+C(d+h+z)/n\) after increasing \(C_0\). The class is bounded, so only
finitely many nonempty shells are needed; the sublevel below \(r_0\) is
absorbed by the remainder. This proves \cref{eq:isomorphic-lemma}. It is the
standard localization-and-peeling scheme~\citep{koltchinskii2006local}, with
the uniform cover making the chaining valid for every empirical measure.
\end{proof}

Applying \cref{lem:local-deviation} to \(f_W=\ell_W-\ell_{W_\lambda^*}\), and using a weighted union bound with failure probability \(e^{-t}\pi_\lambda\), gives simultaneously over \(\lambda\)
\begin{equation}
\cD_w(\wh W_\lambda,W^0)-A_w^2(\lambda)
\le C_{\mathrm{reg},\lambda}
\frac{d_\lambda+\log H_\lambda+x_\lambda+t}{n_1}
+C\varepsilon_{\lambda,\mathrm{opt}}.
\label{eq:within-rate}
\end{equation}
Indeed, insert the approximate ERM inequality into two copies of \cref{eq:isomorphic-lemma} and absorb the quarter-risk terms.

\subsection{Validation selection and proof of the oracle}

\begin{lemma}[Likelihood-ratio Bernstein condition]
\label{lem:lr-bernstein}
If \(u=p_W/q_0\in[e^{-B},e^B]\), then
\[
Q_0^w\left\{\log(q_0/p_W)\right\}^2
\le c_B\cD_w(W,W^0),
\quad
c_B=\sup_{u\in[e^{-B},e^B]}
\frac{(\log u)^2}{u-1-\log u}<\infty.
\]
\end{lemma}

\begin{proof}
Conditional on \((X,E)\), \(\E_{q_0}(u-1)=0\), and hence
\[
\KL(q_0\|p_W)=\E_{q_0}\{u-1-\log u\}.
\]
The scalar ratio defining \(c_B\) extends continuously at \(u=1\), with limit 2, and is bounded on the compact interval. Integrate the scalar inequality and then average over prompts.
\end{proof}

\begin{proof}[Proof of Theorem~\ref{thm:oracle}]
Condition on the training split. Then \(\wh W_\lambda\) is fixed for validation. Bernstein's inequality, Lemma~\ref{lem:lr-bernstein}, and the inequality \(2\sqrt{ab}\le a/4+4b\) imply, with conditional probability at least \(1-e^{-t}\), simultaneously over \(\lambda\),
\begin{equation}
\left|(P_{n_2}-P_0)h_\lambda\right|
\le\tfrac14P_0h_\lambda+C_B\frac{x_\lambda+t}{n_2},
\quad h_\lambda=\log(q_0/p_{\wh W_\lambda}).
\label{eq:validation-deviation}
\end{equation}
Here the weighted union bound uses \(\sum_\lambda\pi_\lambda\le1\). The unknown term \(-\log q_0\) cancels from the empirical comparison in \cref{eq:selector}, so the selector also compares \(P_{n_2}h_\lambda\).

For any comparator \(\lambda\), combine the lower side of \cref{eq:validation-deviation} for \(\wh\lambda\), the selection inequality, and the upper side for \(\lambda\):
\begin{align*}
\tfrac34\cD_w(\wh W,W^0)
&\le P_{n_2}h_{\wh\lambda}
+C_B\frac{x_{\wh\lambda}+t}{n_2}\\
&\le P_{n_2}h_\lambda
+\kappa_B\frac{x_\lambda-x_{\wh\lambda}}{n_2}
+\varepsilon_{\mathrm{sel}}
+C_B\frac{x_{\wh\lambda}+t}{n_2}\\
&\le\tfrac54\cD_w(\wh W_\lambda,W^0)
+C_B\frac{x_\lambda+t}{n_2}
+\kappa_B\frac{x_\lambda-x_{\wh\lambda}}{n_2}
+C_B\frac{x_{\wh\lambda}+t}{n_2}
+\varepsilon_{\mathrm{sel}}.
\end{align*}
Choose \(\kappa_B\ge C_B\) and enlarge constants; the selected-index code term is nonpositive after cancellation. Insert \cref{eq:within-rate}, then take the infimum over \(\lambda\), proving \cref{eq:oracle-hp}.

Under uniform regularity and balanced splits, integrate the exponential tail. Bounded risk controls the small-probability complement and contributes at most \(C/n\) after a standard truncation of the tail integral. This proves \cref{eq:oracle-expectation}.
\end{proof}

\begin{proof}[Proof of Corollary~\ref{cor:same-sample-oracle}]
Subtract the unavailable but candidate-independent empirical teacher log likelihood from the penalized criterion. Apply the assumed isomorphic inequality to the selected candidate and to an arbitrary comparator. A sufficiently large penalty cancels the selected model's complexity; absorb the two quarter-risk terms and optimize over the comparator.
\end{proof}
